\documentclass[11pt,a4paper]{article}

\usepackage[margin=1in]{geometry}
\usepackage{amssymb}
\usepackage{enumitem}
\usepackage{booktabs}
\usepackage{hyperref}
\usepackage{parskip}

\hypersetup{
    colorlinks=true,
    linkcolor=blue,
    urlcolor=blue,
}

\title{\textbf{Math Bootcamp: Content to be Covered}\\[0.4em]
\large For Incoming MSc Economics Students}
\author{Math Camp 2026}
\date{}

\begin{document}

\maketitle

\section*{Overview}

\begin{itemize}[leftmargin=*]
    \item \textbf{Objective:} Bridge the gap between a weaker undergraduate math background and MSc-level economics courses.
    \item \textbf{Duration:} 3 hours per day for 5 days.
    \item \textbf{Method:} Focus on intuition, fundamentals, and confidence-building. Rigorous proofs and advanced techniques are deferred to fall courses; only simple proofs are included where necessary.
    \item \textbf{Emphasis:} Conceptual topics to be covered throughout graduate microeconomics, macroeconomics, econometrics, and mathematical economics.
\end{itemize}

\tableofcontents
\newpage

%==============================================================================
\section{Day 1 --- Foundations (Sections 1.1--1.9)}
%==============================================================================

\textit{Approximate duration: 3 hours. See Day~1 slides (Sections 1.1--1.9).}

\subsection{1.1 --- Notation and Number Systems}
\begin{itemize}[leftmargin=*]
    \item Number systems: $\mathbb{N}$ (counting), $\mathbb{Z}$, $\mathbb{Q}$, $\mathbb{R}$, $\mathbb{R}^n$;
      nesting $\mathbb{N}\subset\mathbb{Z}\subset\mathbb{Q}\subset\mathbb{R}$.
    \item Symbols: $\in$, $\subset$, $\forall$, $\exists$, $\Rightarrow$, $\Leftrightarrow$.
    \item Scalars versus vectors: $\mathbf{x}=(x_1,\ldots,x_n)\in\mathbb{R}^n$;
      dot product $\mathbf{p}\cdot\mathbf{x}=\sum_i p_i x_i$.
    \item Practice: identify the smallest set containing a given object
      (e.g.\ counts in $\mathbb{N}$; prices in $\mathbb{Q}$; order vectors in product sets).
\end{itemize}

\subsection{1.2 --- Logic and Mathematical Statements}
\begin{itemize}[leftmargin=*]
    \item $P$ \textbf{only if} $Q$ ($P\Rightarrow Q$): $Q$ necessary for $P$.
    \item $P$ \textbf{if} $Q$ ($Q\Rightarrow P$): $Q$ sufficient for $P$.
    \item $P$ \textbf{iff} $Q$ ($P\Leftrightarrow Q$): necessary and sufficient.
    \item Contrapositive: $\neg Q\Rightarrow\neg P$ equivalent to $P\Rightarrow Q$.
    \item Running example: IKEA budget set
      $\mathcal{B}=\{\mathbf{x}\in\mathbb{N}^3:\mathbf{p}\cdot\mathbf{x}\le B\}$;
      practice with $\in$, $\subset$, $\forall$, $\exists$, $\Rightarrow$, and contrapositive.
\end{itemize}

\subsection{1.3 --- Sets and Set Operations}
\begin{itemize}[leftmargin=*]
    \item Notation: $\in$, $\notin$, $\subset$, $=$, $\emptyset$, roster and set-builder forms.
    \item Union, intersection, complement, difference; Cartesian product $A\times B$.
    \item Economic sets: budget sets, consumption sets; product spaces such as
      $\mathbb{Q}\times\mathbb{N}$ versus constraint sets like $\mathcal{B}$.
    \item Practice: describe $\mathcal{B}\cap\mathcal{H}$, $\mathcal{B}\cup\mathcal{H}$,
      $\mathcal{B}\setminus\mathcal{H}$; distinguish elements from subsets.
\end{itemize}

\subsection{1.4 --- Functions}
\begin{itemize}[leftmargin=*]
    \item Function $f\colon A\to B$: domain, codomain, range $f(A)\subseteq B$.
    \item Injective (one-to-one) and surjective (onto); codomain versus range.
    \item Operations: $(f\pm g)$, $(cf)$, $(fg)$, composition $(f\circ g)(x)=f(g(x))$.
    \item A rule must assign \textbf{exactly one} output per input (not a correspondence).
    \item Practice: domain, codomain, range, injectivity, and surjectivity
      (e.g.\ partner age as a function of your age).
\end{itemize}

\subsection{1.5 --- Supremum and Infimum}
\begin{itemize}[leftmargin=*]
    \item Example first: reservation price / willingness to pay when the purchase set
      $P=\{p\in\mathbb{R}_+:p<\bar p\}$ has no maximum but $\sup P=\bar p$.
    \item $\max S$ / $\min S$ versus $\sup S$ / $\inf S$ (least upper bound, greatest lower bound).
    \item When extrema are attained versus only approached (e.g.\ $S=[0,1)$, $\sup S=1\notin S$).
    \item If $\max S$ exists, $\max S=\sup S$; always $\inf S\le x\le\sup S$ for $x\in S$.
\end{itemize}

\subsection{1.6 --- Algebra Revision}
\begin{itemize}[leftmargin=*]
    \item Taylor expansion around $a$:
      $f(a+h)=\sum_{k=0}^{\infty}\frac{f^{(k)}(a)}{k!}h^k$ (first-order: $f(a+h)\approx f(a)+f'(a)h$).
    \item Log and exponential rules; series $\ln(1+x)$ and $e^x$.
    \item \textbf{Why logs:} $\Delta\ln x\approx \Delta x/x$ tracks percentage changes;
      elasticity $\varepsilon_{y,x}=\partial\ln y/\partial\ln x$; log--log regression.
    \item Cobb--Douglas: $U=x_1^{\alpha}x_2^{1-\alpha}$, $Y=AK^{\alpha}L^{1-\alpha}$;
      linear in logs; constant elasticities and expenditure shares.
    \item Local log approximation of any smooth $u>0$ near $\bar{\mathbf{x}}$;
      Cobb--Douglas as a modeling choice when used globally.
    \item Homogeneity of degree~$k$: $f(t\mathbf{x})=t^k f(\mathbf{x})$;
      Cobb--Douglas utility and production are degree~$1$ (CRS: double inputs $\Rightarrow$ double output).
\end{itemize}

\subsection{1.7 --- Single-Variable Derivatives}
\begin{itemize}[leftmargin=*]
    \item Formal definition: $f'(a)=\lim_{h\to 0}\frac{f(a+h)-f(a)}{h}$;
      meaning as slope and marginal rate ($MC$, $MU$, $MP_K$).
    \item Rules: constants, sums, power, product, quotient, chain rule;
      $(e^x)'=e^x$, $(\ln x)'=1/x$, $(a^x)'=a^x\ln a$; inverse-function rule (stated).
    \item Practice: cubic cost, monopolist profit, Cobb--Douglas $Y(K)$, composite cost (chain rule).
    \item \textbf{Euler's theorem} (for $C^1$ homogeneous $f$ of degree~$k$):
      $\sum_i x_i f_{x_i}=k f(\mathbf{x})$ (proof via $t\mapsto f(t\mathbf{x})$).
    \item Adding-up problem: under CRS and competitive factor markets,
      $PQ=rK+wL$; IRS ($k>1$) and DRS ($k<1$) and payment exhaustion.
\end{itemize}

\subsection{1.8 --- Multivariate Calculus}
\begin{itemize}[leftmargin=*]
    \item Level curves $\mathcal{L}_c=\{(x,y):f(x,y)=c\}$ and contour maps;
      3D slice picture for $z=f(x,y)$.
    \item Roadmap: partials $\Rightarrow$ total differential $\Rightarrow$ directional derivative
      $\Rightarrow$ cross-partials (Schwarz).
    \item Partial derivatives (limit definition; fix one variable);
      $MP_K=\partial Y/\partial K$ with $L$ fixed.
    \item Total differential: add/subtract trick, MVT on each bracket,
      $dz=f_x\,dx+f_y\,dy$; Cobb--Douglas output example $dY=MP_K\,dK+MP_L\,dL$.
    \item Gradient $\nabla f$: perpendicular to level curves; points uphill;
      $D_{\mathbf{u}}f=\nabla f\cdot\mathbf{u}=\|\nabla f\|\cos\theta$.
    \item Directional derivative along unit direction $\mathbf{u}$;
      example: subsidy requiring input changes in ratio $2{:}1$.
    \item Cross-partials and Schwarz's theorem (stated); $Y_{KL}$ and complements versus substitutes.
\end{itemize}

\subsection{1.9 --- Limits}
\begin{itemize}[leftmargin=*]
    \item Motivation: Berkeley's critique of ``infinitely small'' $h$ in early calculus.
    \item Sequence limit ($\varepsilon$--$N$): quantifier order and convergence.
    \item Function limit ($\varepsilon$--$\delta$): $0<|x-a|<\delta\Rightarrow|f(x)-L|<\varepsilon$;
      example $f(x)=x+1$ for $x\neq 1$, $f(1)=3$, $\lim_{x\to 1}f(x)=2\neq f(1)$.
    \item Continuity at $a$: $\lim_{x\to a}f(x)=f(a)$;
      differentiability at $a$: derivative limit exists.
    \item Differentiable $\Rightarrow$ continuous; converse fails ($|x|$ at $0$; Weierstrass function).
    \item Practice: $\varepsilon$--$\delta$ proofs for $f'(a)=2a$ when $f(x)=x^2$ and for
      $\lim_{x\to 1}f(x)=2$.
\end{itemize}

\subsection{Appendix (Slides)}
\begin{itemize}[leftmargin=*]
    \item Derivations of single-variable rules: constants, sums, powers, product, quotient.
    \item Key limits for $(e^x)'$ and $(\ln x)'$; exponential and logarithm derivative proofs.
\end{itemize}

%==============================================================================
\section{Day 2 --- Calculus II and Linear Algebra: Implicit Functions, Integration, and Matrix Methods}
%==============================================================================

\textit{Approximate duration: 3 hours.}

\subsection{Implicit Function Theorem}
\begin{itemize}[leftmargin=*]
    \item Explicit versus implicit definitions of relationships.
    \item What problem the Implicit Function Theorem solves: local solvability and comparative statics.
    \item \textbf{Core:} statement and use in one variable ($dy/dx=-F_x/F_y$; $dy^*/d\theta=-F_\theta/F_y$).
    \item Sign of $F_y$ and direction of local response.
    \item Application: market equilibrium comparative statics (demand shifter and equilibrium price).
    \item \textbf{Further study (not required):} several variables and matrix/Jacobian form.
\end{itemize}

\subsection{Integration}
\begin{itemize}[leftmargin=*]
    \item Riemann integral: area under a curve.
    \item Fundamental Theorem of Calculus (statement and intuition).
    \item Definite versus indefinite integrals.
    \item Integration as accumulation.
    \item Link to probability: PDF and CDF; CDF as integral of PDF; expected value as integral (conceptual).
    \item Inverse relationship between differentiation and integration (conceptual).
\end{itemize}

\subsection{Leibniz Rule and Interchanging Differentiation and Integration}
\begin{itemize}[leftmargin=*]
    \item Leibniz product rule (review): derivative of a product of functions.
    \item \textbf{Fixed limits:} differentiation under the integral sign when the integrand depends on a parameter but the bounds do not.
    \item \textbf{Variable limits:} when the upper and lower bounds of integration depend on the choice variable (parameter), the derivative has two extra \emph{endpoint} terms from the moving limits, plus an integral of the partial derivative of the integrand.
    \item Informal statement: if $F(\theta) = \int_{a(\theta)}^{b(\theta)} f(x,\theta)\,dx$, then $F'(\theta)$ equals $f(b(\theta),\theta)\,b'(\theta) - f(a(\theta),\theta)\,a'(\theta)$ plus $\int_{a(\theta)}^{b(\theta)} \frac{\partial f}{\partial \theta}(x,\theta)\,dx$ (under regularity conditions).
    \item Economic motivation: consumer surplus, areas under shifting demand curves, and other objects where both the integrand and the integration region change with a parameter.
    \item When we may interchange $\frac{d}{d\theta}$ and $\int$: continuity, differentiability, and boundedness (conceptual); dominated convergence as the measure-theoretic idea behind safe interchange (stated, not proved).
    \item When interchange fails: discontinuities, unbounded integrands, or limits that move irregularly.
    \item Distinction from Fubini (swap order of integrals), MCT, and LDCT (swap limit and integral); preview on Day~3 via uniform convergence.
\end{itemize}

\subsection{Mean Value Theorem}
\begin{itemize}[leftmargin=*]
    \item Geometric intuition: existence of a tangent parallel to a secant.
    \item Rolle's theorem as a special case (statement only).
    \item Economic uses of the Mean Value Theorem: average versus marginal quantities; bounding changes over intervals.
\end{itemize}

\subsection{Taylor Expansion and Linearization}
\begin{itemize}[leftmargin=*]
    \item First-order Taylor approximation.
    \item Linearization of nonlinear functions around a point.
    \item Interpretation as local first-order comparative statics.
    \item When linearization is useful versus when it fails (curvature, nonlinearity).
    \item Preview of second-order Taylor expansion (conceptual only).
\end{itemize}

\subsection{Vectors and Matrices}
\begin{itemize}[leftmargin=*]
    \item Vectors in $\mathbb{R}^n$: notation, components, and geometric interpretation.
    \item Matrices: dimensions, entries, and matrix--vector multiplication.
    \item Basic matrix operations: addition, scalar multiplication, and multiplication.
    \item Transpose of a matrix or vector.
    \item Identity matrix and role in matrix algebra.
\end{itemize}

\subsection{Linear Systems, Inverses, and Determinants}
\begin{itemize}[leftmargin=*]
    \item Systems of linear equations in matrix form ($A\mathbf{x} = \mathbf{b}$).
    \item Inverse of a square matrix: definition and when it exists.
    \item Determinant: definition for $2 \times 2$ and $3 \times 3$ matrices; link to invertibility (conceptual).
    \item Linear independence and dependence of vectors (brief, intuitive treatment).
    \item Rank of a matrix (conceptual: number of linearly independent rows or columns).
\end{itemize}

\subsection{Jacobians and Economic Applications of Linear Algebra}
\begin{itemize}[leftmargin=*]
    \item Jacobian matrix as a collection of partial derivatives.
    \item Dot product and its economic interpretation (e.g., expenditure, inner products in econometrics).
    \item Quadratic forms (brief introduction: notation and positive/negative definiteness, stated not proved).
    \item How matrix notation simplifies multivariate comparative statics;
      full IFT in matrix form (further study).
\end{itemize}

%==============================================================================
\section{Day 3 --- Real Analysis: Topology, Convergence, and Existence}
%==============================================================================

\textit{Approximate duration: 3 hours.}

\subsection{Metric Spaces and Basic Topology in $\mathbb{R}$ and $\mathbb{R}^2$}
\begin{itemize}[leftmargin=*]
    \item Points, distances, and neighborhoods.
    \item Open sets: definition and intuition.
    \item Closed sets: definition and intuition.
    \item Interior, boundary, and closure (conceptual).
    \item Economic relevance: budget sets, feasible sets, production sets; open versus closed constraints.
\end{itemize}

\subsection{Convergence}
\begin{itemize}[leftmargin=*]
    \item Sequences and limits in $\mathbb{R}^n$.
    \item Pointwise convergence of sequences of functions.
    \item Uniform convergence: definition and intuition.
    \item Why uniform convergence matters: exchange of limits; preservation of continuity; interchange of limit and integral; interchange of differentiation and integration (stated, not proved).
\end{itemize}

\subsection{Continuity (Rigorous Intuition)}
\begin{itemize}[leftmargin=*]
    \item $\varepsilon$--$\delta$ definition (light treatment; primarily conceptual).
    \item Continuity on sets versus at a point.
    \item Continuity of demand, supply, profit, and value functions.
    \item Discontinuities: types and economic meaning.
    \item When discontinuities create problems for optimization or equilibrium.
\end{itemize}

\subsection{Compactness}
\begin{itemize}[leftmargin=*]
    \item Definition of compact set (in $\mathbb{R}^n$, via closed and bounded for this level).
    \item Heine--Borel property (statement).
    \item Why compactness matters: existence of maxima and minima for continuous functions.
    \item Weierstrass Extreme Value Theorem (statement).
    \item Maximum and minimum of continuous functions on compact domains.
\end{itemize}

\subsection{Product Spaces}
\begin{itemize}[leftmargin=*]
    \item Cartesian products $\mathbb{R}^n$.
    \item Notation for vectors and product domains.
    \item Constraints and choice sets as subsets of product spaces.
    \item Separability across dimensions (conceptual).
\end{itemize}

\subsection{Separation Theorem}
\begin{itemize}[leftmargin=*]
    \item Convex sets in $\mathbb{R}^n$ (definition sufficient for separation; detailed convexity on Day~4).
    \item Hyperplane separation of disjoint convex sets (picture-based intuition).
    \item Supporting hyperplanes.
    \item Application to welfare economics: Pareto efficiency and competitive equilibrium.
    \item First and Second Welfare Theorems (high-level).
\end{itemize}

\subsection{Fixed Point Theorems}
\begin{itemize}[leftmargin=*]
    \item Fixed point: definition and economic motivation (equilibrium as fixed point).
    \item At least one fixed point theorem (e.g., Brouwer or contraction mapping).
    \item Role in proving existence of equilibrium.
    \item Contraction mapping as a special case (if covered).
\end{itemize}

%==============================================================================
\section{Day 4 --- Optimization}
%==============================================================================

\textit{Approximate duration: 3 hours.}

\subsection{Convex Sets and Convex Functions}
\begin{itemize}[leftmargin=*]
    \item Definition and geometry of convex sets.
    \item Convex and concave functions.
    \item Upper and lower contour sets (level sets).
    \item Why convexity or concavity guarantees global optima from first-order conditions.
\end{itemize}

\subsection{Quasiconcavity and Quasiconvexity}
\begin{itemize}[leftmargin=*]
    \item Definition of quasiconcave and quasiconvex functions via contour sets.
    \item Relationship between concavity/convexity and quasiconcavity/quasiconvexity.
    \item Why quasiconcavity is the standard regularity condition in utility maximization.
    \item When first-order conditions are sufficient for a global optimum under quasiconcavity.
    \item Strict quasiconcavity and uniqueness of optima (conceptual).
    \item Invariance under monotone transformations.
    \item Distinction between quasiconcavity of the objective and convexity of the feasible set in constrained problems.
\end{itemize}

\subsection{Unconstrained Optimization}
\begin{itemize}[leftmargin=*]
    \item Local versus global optima.
    \item First-order conditions: necessary conditions for interior optima.
    \item Second-order conditions: curvature and nature of critical points.
    \item Concavity and convexity of objective functions and uniqueness of optima.
    \item Economic optimization problems as a class (profit, utility, cost minimization).
\end{itemize}

\subsection{Constrained Optimization: Equality Constraints}
\begin{itemize}[leftmargin=*]
    \item Feasible set defined by equality constraints.
    \item Lagrange multiplier method: setup and interpretation.
    \item Lagrange multipliers as shadow prices or marginal values of constraints.
    \item Necessary conditions for constrained optima (Lagrange theorem, stated).
    \item Utility maximization subject to budget constraint (as problem type).
\end{itemize}

\subsection{Constrained Optimization: Inequality Constraints}
\begin{itemize}[leftmargin=*]
    \item Feasible set with inequalities.
    \item Binding versus non-binding constraints.
    \item Complementary slackness (conceptual).
    \item Karush--Kuhn--Tucker (KKT) conditions (statement and interpretation).
    \item Nonnegativity constraints and corner solutions.
\end{itemize}

\subsection{Envelope Theorem}
\begin{itemize}[leftmargin=*]
    \item \textbf{Placement:} after unconstrained and constrained optimization (FOCs, Lagrange, KKT).
    \item Value function of a parameterized problem:
      $V(\theta)=\max_{x} f(x;\theta)$ subject to constraints.
    \item \textbf{Unconstrained envelope theorem:}
      $\frac{dV}{d\theta}=\frac{\partial f}{\partial\theta}\Big|_{x=x^*(\theta)}$
      (no need to differentiate the optimizer $x^*(\theta)$ directly).
    \item \textbf{Constrained envelope theorem:}
      $\frac{dV}{d\theta}=\frac{\partial \mathcal{L}}{\partial\theta}\Big|_{(x^*,\lambda^*)}$
      with binding multipliers capturing constraint movements.
    \item Interpretation: shadow prices and comparative statics from the optimized value.
    \item Economic examples: profit as a function of a cost shifter or output price;
      indirect utility as a function of income and prices;
      expenditure/cost functions and Shephard-type logic (conceptual).
    \item Contrast with naive chain rule through $x^*(\theta)$: envelope theorem drops
      terms that vanish at the optimum.
\end{itemize}

%==============================================================================
\section{Day 5 --- Dynamic Programming}
%==============================================================================

\textit{Approximate duration: 3 hours.}

\subsection{Motivation and Framework}
\begin{itemize}[leftmargin=*]
    \item Sequential decision-making over time.
    \item Intertemporal trade-offs.
    \item Finite versus infinite horizon (finite focus for bootcamp).
    \item Markov structure: current state summarizes relevant history.
\end{itemize}

\subsection{State Variables versus Control Variables}
\begin{itemize}[leftmargin=*]
    \item State: variables that evolve according to law of motion.
    \item Control: variables chosen each period.
    \item Transition equation or law of motion.
    \item Feasibility constraints in dynamic problems.
    \item Classic problem types: consumption--savings, capital accumulation, resource extraction.
\end{itemize}

\subsection{Bellman Equation}
\begin{itemize}[leftmargin=*]
    \item Value function: maximum attainable payoff from a given state.
    \item Principle of optimality (Bellman's idea).
    \item Recursive formulation of dynamic problems.
    \item Bellman equation for finite-horizon problems.
    \item Terminal condition and backward solution logic.
\end{itemize}

\subsection{Backward Induction}
\begin{itemize}[leftmargin=*]
    \item Solving finite-horizon problems from the last period backward.
    \item Relationship between period-by-period optimization and Bellman recursion.
    \item Policy function: mapping from state to optimal control.
\end{itemize}

\subsection{Value Function Iteration (Conceptual)}
\begin{itemize}[leftmargin=*]
    \item Idea of iterating on the Bellman operator.
    \item Convergence to fixed point of Bellman equation (intuition only, no algorithm).
    \item Distinction between solving finite-horizon and infinite-horizon problems.
\end{itemize}

\subsection{Two-Period Models as Building Blocks}
\begin{itemize}[leftmargin=*]
    \item Structure of two-period consumption--savings problems.
    \item Euler equation (conceptual introduction if appropriate).
    \item Link between dynamic programming and standard intertemporal optimization.
\end{itemize}

%==============================================================================
\section{Cross-Cutting Themes (All Days)}
%==============================================================================

\begin{itemize}[leftmargin=*]
    \item Mathematical notation used in economics graduate courses.
    \item Reading and interpreting first-order conditions, second-order conditions, and comparative statics.
    \item Connecting mathematical objects to economic meaning (marginal, shadow price, value function).
    \item Confidence with standard problem templates used in fall courses.
\end{itemize}

%==============================================================================
\section{Suggested Time Allocation}
%==============================================================================

\begin{center}
\begin{tabular}{@{}p{0.12\textwidth}p{0.26\textwidth}p{0.26\textwidth}p{0.26\textwidth}@{}}
\toprule
\textbf{Day} & \textbf{Block 1 ($\sim$1 hr)} & \textbf{Block 2 ($\sim$1 hr)} & \textbf{Block 3 ($\sim$1 hr)} \\
\midrule
1 & 1.1--1.4: notation; logic; sets; functions & 1.5--1.6: sup/inf; algebra; Cobb--Douglas & 1.7--1.9: derivatives; Euler; multivariate; limits \\
2 & Implicit Function Theorem; MVT & Integration; Leibniz rule; Taylor & Vectors; matrices; Jacobians \\
3 & Open and closed sets & Convergence; continuity; product spaces & Compactness; separation; fixed points \\
4 & Convexity; quasiconcavity & Unconstrained optimization & Lagrange; KKT; envelope theorem \\
5 & Framework; state and control & Bellman equation & Backward induction; VFI idea \\
\bottomrule
\end{tabular}
\end{center}

\end{document}
